% ПРИЛОЖЕНИЯ. По А.3: изходен текст на програми, формати на структури от данни,
% екранни форми, получени резултати.
\section*{ПРИЛОЖЕНИЯ}
\label{sec:appendix}
\addcontentsline{toc}{section}{ПРИЛОЖЕНИЯ}

\subsection*{Приложение А. Формат на съобщенията}

Тук се дава пълното описание на видовете съобщения от договора между сървъра и клиента,
описан в Раздел~\ref{sec:design}, заедно с по един пример за всяко съобщение в двата
формата на обмена~\cite{xml10,rfc8259}. \TODO{таблица с полетата на всяко съобщение и
примерите към нея}

\subsection*{Приложение Б. Избрани части от изходния текст}

Тук се включват анализаторът на кадри по WebSocket~\cite{rfc6455}, двата сериализатора и
слоят за свързване на данни на клиента, с коментар на изходния текст.
\TODO{листинги, добавени с \code{lstinputlisting} след като реализацията приключи}

\subsection*{Приложение В. Екранни форми}

Тук се включват екранните снимки на готовия интерфейс: изгледът с числовите показатели,
изгледът с диаграмите и състоянието при прекъсната връзка. \TODO{екранни снимки с
надписи под тях}

\subsection*{Приложение Г. Отчет от измерванията}

Тук се включва суровият отчет от измерванията, върху който стъпват таблиците в
Раздел~\ref{sec:results}, заедно с описанието на средата от Раздел~\ref{sec:method}.
\TODO{отчет от измерванията}
